\input{../tex-inputs/latex-leseansicht-vorspann}

\section[Gustav Schwarzkopf und Paul Marx an Arthur Schnitzler, 13. 8. 1913]{L04255 Gustav Schwarzkopf und Paul Marx an Arthur Schnitzler, 13. 8. 1913}
\nopagebreak\mylabel{L04255v}
\rehead{ }\normalsize\beginnumbering\briefempfaengerindex{Schnitzler, Arthur@\textsc{Schnitzler, Arthur}!zzzMarx, Paul@\emph{von Paul Marx}!1913-08-131@{13. 8. 1913}|(be}\briefempfaengerindex{Schnitzler, Arthur@\textsc{Schnitzler, Arthur}!zzzSchwarzkopf, Gustav@\emph{von Gustav Schwarzkopf}!1913-08-131@{13. 8. 1913}|(be}
\toendnotes[C]{\smallbreak\pagebreak[2]}
\correspDesc{Versand  durch Gustav Schwarzkopf, Paul Marx am 13. 8. 1913 in Wien
\newline{}Übermittlung  am 13. 8. 1913 in Wien
\newline{}Erhalt  durch Arthur Schnitzler im Zeitraum [14. 8. 1913 – 18. 8. 1913?] in Brijuni}\toendnotes[C]{\smallbreak}
\Standort{CUL, Schnitzler, B 96a.}
\physDesc{Postkarte, 538 Zeichen
\newline{}Handschrift Gustav Schwarzkopf: schwarze Tinte, deutsche Kurrent
\newline{}Handschrift Paul Marx: schwarze Tinte, deutsche Kurrent
\newline{}Versand: Stempel: »\nobreak{}\oindex{I., Innere Stadt@\textbf{I., Innere Stadt}, \emph{Verwaltungsgebiet}|pwk}1/1 Wien 8, 13. VIII. 13, 7\nobreak{}«.  }\toendnotes[C]{\smallbreak}\pstart{}{\pb}I. Tiefer Graben 17\oindex{Wien@\textbf{Wien}!I., Innere Stadt@\textbf{I., Innere Stadt}!Tiefer Graben 17@\textbf{Tiefer Graben 17}, \emph{Wohngebäude}|pw}\pend{}\pstart{}II. 6.\pend{}{\bigskip}\pstart{}Herrn\pend{}\pstart{}D\textsuperscript{r} Arthur Schnitzler\pend{}\pstart{}\textsc{Brioni\oindex{Brijuni@\textbf{Brijuni}|pw}}\pend{}\pstart{}\textsc{Istrien\oindex{Istrien@\textbf{Istrien}|pw}}\pend{}{\bigskip}\vspace{1em}
\pstart
           \raggedleft{}{\pb}13. VIII 13\pend
           \vspace{0.5em}
\pstart
           Lieber Arthur,  entſchuldigen Sie, daß ich nicht gleich Ihren Brief
               beantwortet habe; ich hatte i{\geminationm}er noch Beſchwerden und
                  Mahnungen und war{ }ſehr versti{\geminationm}t;
               jetzt{ }ſcheint es nach und nach normal zu werden und geſtern habe ich{ }ſogar einen
               Nachmittagsausflug verſucht, bei dem ich{ }ſelbſtverſtändlich eingeregnet wurde. Ich
               wünſche, daß Ihre neuen Pläne von gutem Wetter begünſtigt werden mögen\pend
           
\pstart
           und grüße Ihre Frau\pwindex{Schnitzler, Olga 17.\,1.\,1882 Wien – 13.\,1.\,1970 Lugano@\textsc{Schnitzler, Olga} (17.\,1.\,1882 Wien – 13.\,1.\,1970 Lugano), \emph{Schauspielerin, Sängerin}|pwv} und Sie herzlichſt.{\\[\baselineskip]}Ihr{\\[\baselineskip]}\spacefill\mbox{Gustav}\pend
           \leftskip=0em{}\selectlanguage{ngerman}\vspace{1em}
\pstart
           \noindent{}{[}hs. Marx:{]} Herzlichſte Grüße Euch Allen\pend
           \pstart \spacefill\mbox{Paul}\pend{}\selectlanguage{ngerman}\endnumbering\briefempfaengerindex{Schnitzler, Arthur@\textsc{Schnitzler, Arthur}!zzzMarx, Paul@\emph{von Paul Marx}!1913-08-131@{13. 8. 1913}|)be}\briefempfaengerindex{Schnitzler, Arthur@\textsc{Schnitzler, Arthur}!zzzSchwarzkopf, Gustav@\emph{von Gustav Schwarzkopf}!1913-08-131@{13. 8. 1913}|)be}\mylabel{L04255h}
\begin{anhang}
\end{anhang}\newcommand{\dateiname}{L04255}\newcommand{\titel}{Gustav Schwarzkopf und Paul Marx an Arthur Schnitzler, 13. 8. 1913}\newcommand{\editorInnen}{Herausgegeben von Jahnke, SelmaMüller, Martin Anton}\input{../tex-inputs/latex-leseansicht-abspann}
